\chapter{Introdução}

Apresente aqui o contexto geral do trabalho: a área de conhecimento em que a pesquisa se insere, os fatos e conceitos que motivam o estudo, e um panorama de como o tema vem sendo tratado. Utilize citações de autores para embasar as afirmações apresentadas \cite{marconi2017metodologia}.

Em seguida, situe o leitor sobre o recorte específico do trabalho: o que exatamente será estudado, investigado ou desenvolvido, e sob qual abordagem (por exemplo, um estudo de caso, uma pesquisa experimental, uma revisão sistemática).

Por fim, descreva brevemente como o trabalho está estruturado, indicando o conteúdo de cada capítulo. Por exemplo:

Este trabalho está organizado em sete capítulos. O Capítulo 2 apresenta o tema e o problema de pesquisa, definindo a questão central a ser investigada. O Capítulo 3 apresenta os objetivos geral e específicos do trabalho. O Capítulo 4 expõe a justificativa da pesquisa, destacando sua relevância. O Capítulo 5 apresenta a fundamentação teórica, abordando os principais conceitos, tecnologias e trabalhos relacionados ao tema. O Capítulo 6 descreve as etapas metodológicas adotadas na pesquisa. Por fim, o Capítulo 7 apresenta o cronograma de execução das atividades previstas.
